\documentclass[11pt]{article}

\usepackage[margin=1in]{geometry}
\usepackage{titling}
\setlength{\droptitle}{-3em}

\title{HW3}
\author{
Greyson Knippel \\[4pt]
CPSC-321- Fall 2026\\[4pt]
}
\date{\today}

\begin{document}

\maketitle

%============================================================
\section*{Query 1}

thresholds: province area $> 100000$, country gdp $> 40000$.

\subsection*{output}
{\small
\begin{verbatim}
 country_code |       country_name       |  gdp  |  province_name   |  area
--------------+--------------------------+-------+------------------+---------
 US           | United States of America | 46900 | Oregon           |  254799
 US           | United States of America | 46900 | Washington       |  184661
 CA           | Canada                   | 43100 | Ontario          | 1076395
 CA           | Canada                   | 43100 | British Columbia |  944735
(4 rows)
\end{verbatim}
}

\subsection*{matching example}
row: US, Oregon.
\begin{verbatim}
country:  ('US', 'United States of America', 46900, 3.8)
province: ('Oregon', 'US', 254799)
\end{verbatim}
oregon's area is 254799 $> 100000$ and the US gdp is 46900 $> 40000$.

\subsection*{nonmatching example}
\begin{verbatim}
country:  ('MX', 'Mexico', 15300, 4.7)
province: ('Sonora', 'MX', 179355)
\end{verbatim}
sonora's area is over the threshold, but mexico's gdp is 15300 $< 40000$. sonora is not in the output.

%============================================================
\section*{Query 2}

thresholds: same as query 1.

\subsection*{output}
{\small
\begin{verbatim}
 country_code |       country_name       |  gdp  |  province_name   |  area
--------------+--------------------------+-------+------------------+---------
 US           | United States of America | 46900 | Oregon           |  254799
 US           | United States of America | 46900 | Washington       |  184661
 CA           | Canada                   | 43100 | Ontario          | 1076395
 CA           | Canada                   | 43100 | British Columbia |  944735
(4 rows)
\end{verbatim}
}

\subsection*{examples}
same as query 1. the output is identical to query 1.

%============================================================
\section*{Query 3}

thresholds: province area $< 100000$, city population between 10000 and 300000 (inclusive).

\subsection*{output}
{\small
\begin{verbatim}
 country_code | province_name | area
--------------+---------------+-------
 DE           | Bavaria       | 70550
 DE           | Saxony        | 18450
(2 rows)
\end{verbatim}
}

\subsection*{matching example}
row: DE, Bavaria.
\begin{verbatim}
province: ('Bavaria', 'DE', 70550)
city:     ('Neustadt', 'Bavaria', 'DE', 13000)
\end{verbatim}
bavaria's area is 70550 $< 100000$ and neustadt's population of 13000 is in the range.

\subsection*{nonmatching example}
\begin{verbatim}
province: ('Jalisco', 'MX', 78599)
city:     ('Guadalajara', 'Jalisco', 'MX', 1385629)
city:     ('Zapopan', 'Jalisco', 'MX', 1476491)
city:     ('Tlaquepaque', 'Jalisco', 'MX', 687127)
\end{verbatim}
jalisco's area is under the threshold, but every city is over 300000. jalisco is not in the output.

%============================================================
\section*{Query 4}

thresholds: same as query 3.

\subsection*{output}
{\small
\begin{verbatim}
 country_code | province_name | area
--------------+---------------+-------
 DE           | Bavaria       | 70550
 DE           | Saxony        | 18450
(2 rows)
\end{verbatim}
}

\subsection*{examples}
same as query 3. the output is the same as query 3.

%============================================================
\section*{Query 5}

thresholds: city population $> 500000$, country inflation $< 4.0$.

\subsection*{output}
{\small
\begin{verbatim}
 country_code | country_name | inflation | province_name |  area
--------------+--------------+-----------+---------------+---------
 CA           | Canada       |       2.4 | Ontario       | 1076395
 DE           | Germany      |       2.2 | Bavaria       |   70550
 DE           | Germany      |       2.2 | Saxony        |   18450
(3 rows)
\end{verbatim}
}

\subsection*{matching example}
row: CA, Ontario.
\begin{verbatim}
country:  ('CA', 'Canada', 43100, 2.4)
province: ('Ontario', 'CA', 1076395)
city:     ('Toronto', 'Ontario', 'CA', 2794356)
city:     ('Ottawa', 'Ontario', 'CA', 1017449)
\end{verbatim}
ontario has two different cities over 500000 and canada's inflation is 2.4 $< 4.0$.

\subsection*{nonmatching example}
\begin{verbatim}
country:  ('US', 'United States of America', 46900, 3.8)
province: ('Washington', 'US', 184661)
city:     ('Seattle', 'Washington', 'US', 737015)
city:     ('Spokane', 'Washington', 'US', 228989)
city:     ('Vancouver', 'Washington', 'US', 190915)
\end{verbatim}
the US inflation is under the threshold, but only seattle is over 500000. washington is not in the output.

%============================================================
\section*{Query 6}

thresholds: same as query 5.

\subsection*{output}
{\small
\begin{verbatim}
 country_code | country_name | inflation | province_name |  area
--------------+--------------+-----------+---------------+---------
 CA           | Canada       |       2.4 | Ontario       | 1076395
 DE           | Germany      |       2.2 | Bavaria       |   70550
 DE           | Germany      |       2.2 | Saxony        |   18450
(3 rows)
\end{verbatim}
}

\subsection*{examples}
same as query 5. the output is the same as query 5.

%============================================================
\section*{Query 7}

no thresholds.

\subsection*{output}
{\tiny
\begin{verbatim}
 city_name_1 | province_name_1 | country_code_1 | population_1 | city_name_2 | province_name_2  | country_code_2 | population_2
-------------+-----------------+----------------+--------------+-------------+------------------+----------------+--------------
 Vancouver   | Washington      | US             |       190915 | Vancouver   | British Columbia | CA             |       662248
 Portland    | Ontario         | CA             |        15000 | Portland    | Oregon           | US             |       652503
(2 rows)
\end{verbatim}
}

\subsection*{matching example}
row: Vancouver (US), Vancouver (CA).
\begin{verbatim}
city: ('Vancouver', 'Washington', 'US', 190915)
city: ('Vancouver', 'British Columbia', 'CA', 662248)
\end{verbatim}
same name, different countries, different populations. the less populous city (US) is first.

\subsection*{nonmatching example}
\begin{verbatim}
city: ('Neustadt', 'Bavaria', 'DE', 13000)
city: ('Neustadt', 'Saxony', 'DE', 12000)
\end{verbatim}
same name, but both are in DE. neustadt is not in the output.

%============================================================
\section*{Query 8}

threshold: border length $> 400$.

\subsection*{output}
{\small
\begin{verbatim}
 country_code |       country_name
--------------+--------------------------
 CA           | Canada
 US           | United States of America
(2 rows)
\end{verbatim}
}

\subsection*{matching example}
row: US.
\begin{verbatim}
country: ('US', 'United States of America', 46900, 3.8)
country: ('MX', 'Mexico', 15300, 4.7)
border:  ('MX', 'US', 3145)
\end{verbatim}
the US has a higher gdp (46900 $>$ 15300) and lower inflation (3.8 $<$ 4.7) than mexico, and the border is 3145 $> 400$. the US is in \texttt{country\_code\_2}.

\subsection*{nonmatching example}
\begin{verbatim}
country: ('MX', 'Mexico', 15300, 4.7)
border:  ('MX', 'US', 3145)
border:  ('CA', 'MX', 500)
\end{verbatim}
mexico has a lower gdp and higher inflation than both of its neighbors. mexico is not in the output.

%============================================================
\section*{Query 9}

threshold: same as query 8.

\subsection*{output}
{\small
\begin{verbatim}
 country_code |       country_name
--------------+--------------------------
 US           | United States of America
 CA           | Canada
(2 rows)
\end{verbatim}
}

\subsection*{examples}
same as query 8. the output has the same rows as query 8 (order differs, no sorting required).

%============================================================
\section*{Query 10}

question: which provinces have a city with over 1000000 people, in a country with inflation under 4.0? \\
thresholds: city population $> 1000000$, country inflation $< 4.0$. \\
sorted by province area, largest first, so the largest regions are listed first.

\subsection*{output}
{\small
\begin{verbatim}
 country_name | province_name |  area
--------------+---------------+---------
 Canada       | Ontario       | 1076395
 Germany      | Bavaria       |   70550
(2 rows)
\end{verbatim}
}

\subsection*{matching example}
row: Germany, Bavaria.
\begin{verbatim}
country:  ('DE', 'Germany', 52100, 2.2)
province: ('Bavaria', 'DE', 70550)
city:     ('Munich', 'Bavaria', 'DE', 1512491)
\end{verbatim}
munich is over 1000000 and germany's inflation is 2.2 $< 4.0$.

\subsection*{nonmatching example}
\begin{verbatim}
country:  ('MX', 'Mexico', 15300, 4.7)
province: ('Jalisco', 'MX', 78599)
city:     ('Guadalajara', 'Jalisco', 'MX', 1385629)
city:     ('Zapopan', 'Jalisco', 'MX', 1476491)
\end{verbatim}
jalisco has cities over 1000000, but mexico's inflation is 4.7 $> 4.0$. jalisco is not in the output.

\end{document}
